\section{Evaluation}
\label{sec:evaluation}

We evaluate CCV along four research questions:

\begin{itemize}
    \item \textbf{RQ1}: How often do LLM agents hallucinate about code?
    \item \textbf{RQ2}: Does deterministic verification catch hallucinations that LLM-as-judge misses?
    \item \textbf{RQ3}: Does claim-calibrated confidence improve downstream decision quality?
    \item \textbf{RQ4}: Which claim types are most commonly hallucinated?
\end{itemize}

\subsection{Dataset}

We evaluate on the CyberGym benchmark~(ICLR 2026): 201 real-world C/C++ vulnerability repositories spanning 188 projects including binutils, curl, ffmpeg, ghostscript, libxml2, openssl, php, wireshark, and yara. Each repository contains the vulnerable source code and a vulnerability description from the CyberGym HuggingFace dataset. We generate LLM reasoning about each vulnerability using 7 models, then run CCV to extract and verify claims against the actual source code.

\subsection{Models}

We evaluate 7 models across commercial and open-source tiers (Table~\ref{tab:models}), accessed via Google Vertex AI and Red Hat Models.corp.

\begin{table}[t]
\centering
\caption{Models evaluated. Open models accessed via Red Hat Models.corp (OpenAI-compatible API on internal vLLM infrastructure).}
\label{tab:models}
\small
\begin{tabular}{llrl}
\toprule
\textbf{Model} & \textbf{Provider} & \textbf{Size} & \textbf{Type} \\
\midrule
Claude Sonnet 4 & Anthropic (Vertex) & Very Large & Commercial \\
Claude Haiku 4.5 & Anthropic (Vertex) & Very Large & Commercial \\
Llama 3.3 70B FP8 & Red Hat AI & 70B & Open \\
GPT-OSS 20B & OpenAI (open) & 20B & Open \\
Qwen3 14B & Qwen & 14B & Open \\
Granite 3.3 8B & IBM & 8B & Open \\
Mistral 7B v0.3 & Mistral & 7B & Open \\
\bottomrule
\end{tabular}
\end{table}

Claim extraction uses Claude Sonnet 4 for all samples (isolating extraction quality from reasoning quality).

\subsection{Results}

\subsubsection{RQ1: Hallucination Frequency}

Across 1{,}278 vulnerability analyses (201 repos $\times$ up to 7 models), LLMs produced 23{,}251 factual claims about code (Table~\ref{tab:rq1}). CCV refuted 6{,}839 of these (31.1\% aggregate hallucination rate).

\begin{table}[t]
\centering
\caption{Hallucination rate by model on 201 CyberGym repos. Sorted by rate. 23{,}251 total claims verified.}
\label{tab:rq1}
\small
\begin{tabular}{lrrrrr}
\toprule
\textbf{Model} & \textbf{Repos} & \textbf{Claims} & \textbf{Verified} & \textbf{Refuted} & \textbf{Halluc.\ \%} \\
\midrule
Granite 3.3 8B & 201 & 3{,}272 & 2{,}232 & 552 & 16.9\% \\
Llama 3.3 70B & 201 & 4{,}944 & 3{,}780 & 1{,}032 & 20.9\% \\
Mistral 7B & 201 & 2{,}382 & 1{,}575 & 606 & 25.4\% \\
GPT-OSS 20B & 110 & 2{,}192 & 1{,}427 & 681 & 31.1\% \\
Qwen3 14B & 163 & 2{,}401 & 1{,}436 & 791 & 32.9\% \\
Claude Haiku 4.5 & 201 & 4{,}205 & 2{,}452 & 1{,}639 & 39.0\% \\
Claude Sonnet 4 & 201 & 3{,}855 & 2{,}237 & 1{,}538 & 39.9\% \\
\midrule
\textbf{All models} & --- & \textbf{23{,}251} & \textbf{15{,}139} & \textbf{6{,}839} & \textbf{31.1\%} \\
\bottomrule
\end{tabular}
\end{table}

The most striking finding is that open models outperform commercial models on factual accuracy about code. Granite 3.3 8B (16.9\%) and Llama 3.3 70B (20.9\%) produce fewer hallucinated claims than Claude Sonnet 4 (39.9\%) and Claude Haiku 4.5 (39.0\%). This reversal of the expected capability ordering is driven by claim type distribution: Claude models produce more \textsc{Absence} and \textsc{File\_Classification} claims, which have higher refutation rates. Open models produce more \textsc{File\_Exists} and \textsc{Function\_Exists} claims, which are easier to verify correctly.

Model size does not straightforwardly predict accuracy within the open model tier: the 8B Granite (16.9\%) outperforms the 70B Llama (20.9\%), the 14B Qwen (32.9\%), and the 20B GPT-OSS (31.1\%).

\subsubsection{RQ2: CCV vs.\ LLM-as-Judge}

CCV's zero false-\textsc{Verified} rate is its key safety property: when CCV says a claim is correct, it \emph{is} correct. This guarantee holds by construction because the verification tools (grep, \texttt{os.path.isfile}, lockfile parsing) are deterministic. In our evaluation on 23{,}251 claims, CCV produced zero false-\textsc{Verified} errors across all 7 models.

In the LLM-as-judge comparison on a subset ($n=263$ disagreements), CCV was correct 97\% of the time (255 vs.\ 8). The judge's primary failure mode is confirming hallucinated claims that CCV correctly refutes.

\subsubsection{RQ3: Downstream Impact}

CCV's action distribution demonstrates practical impact: across all analyses, only 16--26\% receive \textsc{Boost} (all claims verified). The majority receive \textsc{Flag} (mixed results, human review needed) or \textsc{Override} (majority hallucinated, don't trust). Without CCV, all analyses would have been trusted at face value.

\subsubsection{RQ4: Per-Type Analysis}

Hallucination rates vary significantly by claim type. \textsc{File\_Exists} claims are nearly always correct (verified via \texttt{os.path.isfile}). \textsc{Absence} claims (``no sanitization exists,'' ``no callers found'') are the most commonly hallucinated, consistent with the known difficulty of proving negatives through LLM reasoning. \textsc{Function\_Exists} and \textsc{Import\_Exists} claims fall in between, with accuracy depending on how precisely the LLM specifies file paths and function names.

The aggregate ECE across controlled fixtures is 0.167, indicating that CCV's per-type confidence values are systematically conservative (actual accuracy exceeds predicted confidence). This is a desirable property for a safety-critical tool.
